% chapters/ch20_strong_engine.tex
% Chapter 20: Strong Engine Architecture — Caged Beast Model
% Source: NRMO_v5_updated.pdf, Chapter 20 (lines 9147-9322 of v5_layout.txt)

\chapter{Strong Engine Architecture --- Caged Beast Model}
\label{ch:strong-engine}

\nrmobox{Document status}{
\begin{tabular}{ll}
\textbf{Status}: & Structural Requirement (20-World Consensus) \\
\textbf{Source}: & 10 experts $\times$ approx.~300 years $\times$
20 world lines \\
\textbf{Classification}: & Core Architecture Extension
\end{tabular}\\[0.4em]
NRMO is a constitution to prevent breaking; a separate strong engine
carries the force to advance the world. Both are separated; NRMO
only vetoes breakdown.
}

% =====================================================================
\section{v7 Realignment: Real-Side Primacy and Two-Layer Structure}
\label{sec:strong-engine-v7-realignment}

\nrmobox{v7 structural correction}{
Earlier versions misplaced the Strong Engine $\Omega$ Full on the
civ-sim side. The corrected architecture places it on the \textbf{real
side} as the primary actor. The civ-sim side hosts a separate
MaxForwardEngine that imitates the Strong Engine.
}

\paragraph{Real side --- Strong Engine $\Omega$ Full.}
The Strong Engine performs all exploration, reasoning, and decision-making
(unrestricted scope) over the \emph{true dynamics of the target domain}.
Its default is \textbf{maximum forward movement}. NRMO removes only the
components heading toward absorbing ruin; the magnitude of forward movement
itself is never reduced. Stagnation and ``minimum viable'' choices are
treated as a form of ruin (passive death) and excluded. The Loom layer
resides on this real side as the ruin-boundary governance.

\paragraph{civ-sim side --- MaxForwardEngine.}
A separate MaxForwardEngine on the civilisation-simulation side imitates
the Strong Engine. It tolerates aggressive death/ruin (for feedback),
excludes extreme caution, and computes the attack-peak and defense-peak,
seeking the best score. The best action found is returned to the real-side
Strong Engine, which advances without breaking. The two sides exchange and
accumulate information bidirectionally.

\paragraph{World-dependence.}
Maximum forward is the \emph{default} posture, not a universal optimum.
The civ-sim layer explores each world's true optimum: aggression dominates
in some worlds (inventory-cycle domains), sustainable growth in others
(long-horizon civilisations), and extreme caution in yet others
(V9-minimum-type worlds). The two-layer structure \emph{discovers}, rather
than imposes, the optimum.

\paragraph{Rollout discipline (root-cause fix).}
The engine's rollout must use the target domain's true dynamics. Embedding
a fixed civ-sim model inside the engine causes divergence and over-defensive
collapse --- the root cause of prior ``poor numbers.'' A short evaluation
horizon compounds this by failing to credit long-run compounding growth,
biasing the choice toward over-defense. See the Realignment Master Spec for
the full component mapping.

% =====================================================================
\section{Core Consensus}
\label{sec:strong-engine-consensus}

\subsection{Consensus 1: NRMO Alone Does Not Advance}

\paragraph{Structural limitation.}
NRMO is extremely strong at ``not breaking.'' However, the following
require propulsive force (engine):

\begin{itemize}
\item Technological breakthroughs.
\item Non-continuous market jumps.
\item Social / relational phase transitions.
\end{itemize}

\paragraph{Formal proof: limitation of NRMO alone.}

\begin{itemize}
\item \textbf{Premise 1}: NRMO constraint $= P(\mathrm{Ruin} \mid a)
\le \varepsilon$.
\item \textbf{Premise 2}: $\text{threshold breakthrough} \subseteq
\{a \mid P(\mathrm{Ruin} \mid a) > \varepsilon\}$.
\item \textbf{Therefore}: NRMO alone cannot break through threshold
structurally.
\end{itemize}

\begin{equation}
\forall w \in \mathrm{Worlds}:\;
\mathrm{NRMO\_only}(w) \;\Rightarrow\;
\neg \mathrm{CanCrossThreshold}(w).
\label{eq:nrmo-alone-limitation}
\end{equation}

\subsection{Consensus 2: A Strong Engine Is Essential}

20 world-line observations:

\begin{itemize}
\item Only weak repetitive engine $\to$ survives but loses
initiative.
\item Lacking strong engine $\to$ eliminated in discontinuous
environments.
\end{itemize}

\paragraph{Conclusion.}
The theory that an engine is unnecessary is a denial of reality.

% =====================================================================
\section{Rejected Designs (Commonalities Across Failed Worlds)}
\label{sec:strong-engine-failures}

Three patterns that failed in all world lines:

\subsection{Failure 1: Engine Placed Inside NRMO}

Placing the strong engine under NRMO constraints makes
non-continuous jumps impossible.

\begin{equation}
\mathrm{Engine} \subseteq \mathrm{NRMO}
\;\Rightarrow\;
\text{All actions subject to } P(\mathrm{Ruin}) \le \varepsilon.
\end{equation}

But strong engine requires $P(\mathrm{Ruin}) > \varepsilon$ for
non-continuous jumps $\Rightarrow$ contradiction $\Rightarrow$
failure.

\subsection{Failure 2: NRMO Intervening in Engine Purpose / Strategy}

When NRMO says what ``should be done,'' the constitution becomes a
ruler.

\subsection{Failure 3: Suppressing Propulsive Force Under ``Safety''}

Ruin avoidance becomes an end in itself, prohibiting exploration
altogether.

\paragraph{Common consequence.}

\begin{itemize}
\item Religionisation (NRMO absolutised).
\item Stagnation (disappearance of exploration).
\item Self-exculpation (``safety'' as excuse).
\item Disappearance of decision subject (outsourcing of
responsibility).
\end{itemize}

% =====================================================================
\section{Correct Structure (Design That Survived Across Worlds)}
\label{sec:strong-engine-correct}

\begin{table}[ht]
\centering
\small
\begin{tabularx}{\linewidth}{lXX}
\toprule
\textbf{Layer} & \textbf{Role} & \textbf{Authority} \\
\midrule

NRMO (Constitution)
& Veto of irreversible breakdown
& Veto only (NO / HOLD / SAFE\_EXIT) \\

Engine (Propulsion)
& Non-continuous jump / threshold breakthrough
& Exploration, experiment, challenge \\

Prohibition
& Engine: ignoring NRMO; NRMO: propulsion / optimisation / success
definition
& --- \\

\bottomrule
\end{tabularx}
\caption[Caged beast structure]{Correct two-layer architecture
surviving across all world lines.}
\label{tab:caged-beast-structure}
\end{table}

\subsection{Definition of Strong Engine}

\begin{itemize}
\item \textbf{High-speed}: fast repetition rate.
\item \textbf{Single-purpose orientation}: has clear goals.
\item \textbf{Non-continuous jumps}: force to cross thresholds.
\item \textbf{Failure-premise}: many failures built in.
\item \textbf{World-bending}: has force to change the environment.
\end{itemize}

% =====================================================================
\section{Formal Definition of Complete System}
\label{sec:strong-engine-formal}

\begin{equation}
\mathrm{System} = (\mathrm{NRMO},\, \mathrm{Engine},\,
\mathrm{Separation}).
\end{equation}

\begin{equation}
\mathrm{Separation} := \begin{cases}
\mathrm{Engine} \not\subseteq \mathrm{NRMO}, \\
\mathrm{NRMO.authority} = \mathrm{veto\_only}, \\
\mathrm{Engine} \to \mathrm{NRMO}: \text{cannot ignore}.
\end{cases}
\label{eq:separation-definition}
\end{equation}

\paragraph{Health metric.}

\begin{equation}
\mathrm{health}(\mathrm{System}) \propto
\frac{1}{\mathrm{veto\_frequency}}.
\label{eq:health-metric}
\end{equation}

% =====================================================================
\section{Operation Flow (Integrated)}
\label{sec:strong-engine-flow}

\begin{lstlisting}[basicstyle=\ttfamily\small,frame=single]
Engine: EXPLORE -> ATTEMPT -> FAIL/SUCCEED -> ITERATE
       |
       v (each action)
NRMO:  CHECK P(Ruin | a) <= epsilon ?
       |-- YES: PASS  (Engine continues)
       |-- NO:  VETO  (HOLD / SAFE_EXIT)

Engine.purpose   <-- NRMO has no access
Engine.strategy  <-- NRMO has no access
Engine.success   <-- NRMO has no access
\end{lstlisting}

% =====================================================================
\section{Cross-Domain Successful Implementation Examples}
\label{sec:strong-engine-examples}

\begin{table}[ht]
\centering
\small
\begin{tabularx}{\linewidth}{lXX}
\toprule
\textbf{Domain} & \textbf{NRMO role} & \textbf{Engine role} \\
\midrule

Startup
& Cashflow runway management
& Product-market fit search \\

Investment
& Risk-Off trigger / max drawdown cap
& Alpha generation \\

Relationship
& Irreversibility boundary monitoring
& Commitment / deepening \\

Career
& Reputation damage prevention
& Skill acquisition / challenge \\

\bottomrule
\end{tabularx}
\caption[Strong engine cross-domain examples]{Cross-domain successful
implementation examples of the caged beast architecture.}
\label{tab:strong-engine-examples}
\end{table}

% =====================================================================
\section{Health Metrics}
\label{sec:strong-engine-health}

NRMO health is measured by \emph{decrease} in intervention
frequency.

\begin{itemize}
\item \textbf{Healthy}: veto frequency decreasing
(Engine learning NRMO constraints).
\item \textbf{Warning}: veto frequency rising (Engine drift or NRMO
over-sensitive).
\item \textbf{Pathological}: veto frequency very high (NRMO becoming
religion).
\end{itemize}

% =====================================================================
\section{Implementation Checklist}
\label{sec:strong-engine-checklist}

\begin{itemize}
\item[$\square$] Is NRMO a veto-only constitution?
\item[$\square$] Is Engine outside NRMO?
\item[$\square$] Does NRMO not define ``what should be done''?
\item[$\square$] Does NRMO not define ``how to succeed''?
\item[$\square$] Is NRMO exercising only veto authority?
\item[$\square$] Is intervention frequency decreasing over time?
\end{itemize}

% =====================================================================
\section{Final Declaration}
\label{sec:strong-engine-final}

\nrmobox{Caged Beast Model}{
\centering
\textit{NRMO without an Engine stays still and eventually fails.}\\
\textit{An Engine without NRMO breaks and eventually fails.}\\[0.3em]
\textit{The caged beast model is the only architecture that survived
in all world lines.}
}

% =====================================================================
\section*{Connections to Other Parts}

\begin{itemize}
\item Equations~\ref{eq:separation-definition} and
\ref{eq:health-metric} are the formal core of the
governance--execution separation invariant
(Invariant~\ref{inv:gov-exec-sep}) made operational.
\item Cross-domain examples find their engineering instantiation in
$\Omega$ Full (Part~IX): the engine is the StrongEngine $\Omega$
Full, the constitution is the NRMO Core veto layer, and the
formal separation is enforced architecturally in the codebase
(Part~X).
\item The 20-world consensus result is the empirical foundation for
the Part~XI validation of $\Omega$ Full v5.0 across the five
world families.
\end{itemize}

% =====================================================================
\section{v7 Phase E: DecisionCompass Two-Layer with Real Governance}
\label{sec:strong-engine-v7-phase-e}

\nrmobox{v7 Phase E (2026-05-30)}{
DecisionCompass is realigned as a fully real-side personal-decision application:
the engine is the Strong Engine $\Omega$ Full (with an internal MaxForward
civ-sim), the constitutional boundary is the real \texttt{nrmo\_core} veto, and
the dynamics are the real \texttt{civstate\_transition}. The user is shown only
the ``non-ruin maximum forward'' recommendation.
}

The governance/execution separation is preserved structurally: governance builds
the admissible set $A_t = \mathrm{NRMO}(X_t)$ using veto only (yes/no/hold), and
the engine searches strictly within it, $a_t = \mathrm{Engine}(A_t)$. The engine
never sees the veto thresholds and cannot override the admissible set.

\subsection{Validation (reproducible)}
\begin{itemize}
  \item \textbf{Admissible set shrinks with risk.} Healthy states admit a
        forward sub-spectrum; an overheated state admits fewer; an edge state
        (high exposure, weak governance) admits none and returns \textsc{exit\_hold},
        i.e.\ change direction (lower exposure) before advancing.
  \item \textbf{Separation enforced.} The engine's chosen action always lies in
        the admissible set; vetoed aggressive candidates are unreachable.
  \item \textbf{Trajectory.} Over 200 personal-decision steps the system records
        ruin 0\% while continuing to advance.
\end{itemize}

\nrmobox{Honest note}{
In this personal domain the recommended forward level is conservative because
higher growth raises exposure $X$ and lowers the prosperity measure
$R+E+G+O+K-X$; among equally surviving admissible actions the lower-growth one
wins by a small margin. This is world-dependent caution (it is still forward
motion, not stagnation), and the aggressive options were removed by governance,
not by the engine. The magnitude depends on the exposure weight, which matches
the existing \texttt{CivilizationDynamics} convention.
}
